%% ----------------------------------------------------------------
%% Chapter 2 -- Literature Survey
%% ----------------------------------------------------------------
\chapter{Literature Survey}
\label{chap:litsurvey}

\section{Review of Previous Works}
\label{sec:priorwork}

Prior art exists across five different areas; ISAC network topology, empirical characterization of Wi-Fi FTM performance, mmWave radar capability analysis, covariance intersection theory application, and age of information modeling. With respect to both ISAC and 6G; researchers such as Wang et al.~\cite{wang_isac} and ETSI GR ISC 001~\cite{etsi_isc_001} have shown that there needs to be a combination of millimeter-wave band and sub-6~GHz for robust tracking of multiple targets. Therefore our dual modality system using both 2.4~GHz and 24~GHz bands is validated.

In regards to Multistatic Radar, Pegoraro et al.~\cite{pegoraro_asymov} explored an 802.11ay 60~GHz multistatic ISAC which illustrated how spatial diversity can overcome occlusions. In contrast to a multiple receiver based methodology, we will be utilizing a singular node as a fallback.

For Wi-Fi FTM, Aggarwal et al.~\cite{aggarwal_wifi} demonstrated meter-level LOS (line-of-sight) accuracy at frequencies ranging from 80 to 160~MHz. However they also identified a degradation of 5 meters due to multipath. This validates both the stochastic limits for FTM and therefore supports the utilization of radar for near field calibration.

Asynchronous Covariance Intersection (Chen~\cite{chen_ci} \& Yang~\cite{yang_heterogeneous}) established that ACI guarantees non-divergent fusion even in the case where all cross correlations between the different sources are unknown, while extending this result to cover heterogeneous track fusion. Therefore, ACI will provide the linear algebra foundation upon which our implementation of DCI will rely.

Age of information (Noroozi~\cite{noroozi_aoi}) provided upper bounds on deviation of information for variance inflation caused by communication delays. These results form the basis for the dynamic coast penalty terms we have proposed.

\textbf{ISAC in 6G.} Wang~\cite{wang_isac} established a taxonomy for ISAC into two main categories: Communication Centric and Sensing Centric Designs. ETSI~\cite{etsi_isc_001} enumerated 18 ISAC use cases and concluded that single frequency operations would not be sufficient - concurrent sensing with cross source fusion will be necessary for successful deployment of 6G networks. The Next G Alliance Phase~I report~\cite{nextg_phase1} reinforced the economic viability of ISAC as a function of fusing network sensed data with embedded sensor input from devices.

\textbf{Millimetre Wave Sensing.} Pegoraro et al.~\cite{pegoraro_asymov} demonstrated that using an ISAC system at 60~GHz was able to capture the micro-doppler signature of a target in granular detail; however, if the target was out of line-of-sight, then the target would be completely lost as there was no diffraction effect. Pegoraro et al.'s~\cite{pegoraro_asymov} solution - through use of spatial diversity achieved by the placement of distributed receivers at different locations on the network and timing compensation - was able to provide simultaneous tracks, however this required a number of nodes across the network. Monostatic architectures are still vulnerable to losing track without an additional sensing modality.

\textbf{Wi-Fi FTM.} Aggarwal~\cite{aggarwal_wifi} characterized the empirical performance of IEEE 802.11mc FTM ranging. It was possible to achieve accuracies of 1--2 metres when the measurement was made in line-of-sight; however, if measurements were made over paths with significant levels of NLOS multipath, then the errors increased significantly to around 5 metres with a marked positive bias resulting from architectural reflections. The FTM ranging error profile had a high level of spatial variance and exhibited receiver saturation effects close to the transmitter. Machine learning classification techniques were able to classify propagation conditions; however, they did not reduce the inherent stochastic nature of the ranging process. FTM provided long-range identity linked tracking, but did not provide sufficient data to support sub-metre localisation on its own.

\textbf{Covariance Intersection.} Chen~\cite{chen_ci} provided an analysis of CI as a viable sensor data fusion method that consistently fuses data from sensors whose relative motion and/or unknown inter-sensor cross-correlation exist, and developed equations for finding the optimal linear combination of sensors that minimises the upper bound on the fused estimation error. Yang~\cite{yang_heterogeneous} expanded upon this work by applying CI to asynchronous heterogeneous sensor data fusion. Specifically he addressed the issue of aligning the differing dimensionalities of range (a 1-dimensional value) and position (a 2-dimensional vector) measurement spaces by projecting all measurements into a common time reference frame prior to performing the CI calculation.

\textbf{Age of information.} Noroozi~\cite{noroozi_aoi} established age of information (AoI) as a measure of data freshness in CPS's. He showed mathematically that Periodic Sampling is Suboptimal for Intermittent/Event Driven Updates such as Bursty First Target Measurement (FTM)~\cite{noroozi_aoi}. The DoI framework gives upper-bounds to degrade the quality of estimates with respect to how much time has passed since the most recent valid measurement. This provides the theoretical foundation for the use of covariance inflation in order to force all filter engines to de-weight old measurements rather than place implicit trust in them.

\section{Gaps Identified in Prior Research Approaches}
\label{sec:gaps}

Five voids appear in literature.

\textbf{Gap 1: Freshness Blindness.} Current Filter Methods treat stale WiFi readings as permanently valid between 1~Hz updates. This causes Trajectory Overshoot when the target performs maneuvers. While AoI variance inflation theory does exist~\cite{noroozi_aoi}, it is rarely implemented in real-time arithmetic for low-cost microcontrollers.

\textbf{Gap 2: Hardware Cost Barrier.} Most current prototypes of ISACs require expensive, proprietary hardware. There is a lack of accessible multi-sensor architecture capable of performing heterogeneous fusion using off-the-shelf commercial equipment.

\textbf{Gap 3: Single-source occlusion vulnerability.} Unimodal radar cannot verify the identities of targets and will completely fail if blocked. Solutions utilizing multistatic spatial diversity require large amounts of distributed infrastructure~\cite{pegoraro_asymov}. The lack of understanding using a dynamic mathematical backup from a single node as a secondary, penetrating method such as Wi-Fi FTM is uncharted.

\textbf{Gap 4: The static nature of covariance rigidity.} Current static implementations~\cite{chen_ci,yang_heterogeneous} utilize constant variance matrices developed for a uniform environment and do not account for the inverse of distance-dependent errors in Wi-Fi (i.e., high near-field error rates, higher accuracy at longer ranges) and radar (highly accurate at short distances, degrades rapidly with increasing distance). Utilizing static confidence ratios will lead to less than optimal performance for tracking objects across different range bands.

\textbf{Gap 5: Edge machine learning precision limitations.} INT8 quantization - which is standard for edge deployments - converts continuous values to discrete representations (steps) causing the loss of sub centimeter regression accuracy. Optimized float32 versions of machine learning algorithms that handle asynchronous time-series data without naive zero-filling strategies have limited availability.

In summary, the five areas that need to be addressed are as follows: Real-time AoI covariance inflation is needed to address freshness blindness. Hardware exclusivity will be addressed via implementation of an operating system capable of running on two cores and commercially available. Occlusion vulnerability will be addressed through implementing 2.4~GHz non line-of-sight as a dynamically implemented fallback. Dynamic or learned via ML trust weighting will need to be utilized to address static covariance. Last, the precision limitations of edge devices will need to be addressed by implementing a Float32 MLP with forward fill-based temporal encoding and freshness awareness of its input features. This project addresses these gaps by developing real-time AoI driven covariance inflation as part of Iteration~1 and a precision preserving Float32 MLP with forward fill-based temporal encoding as part of Iteration~2 on a low cost hardware pipeline.
